\section{System Design}
\label{sec:system-design}

This section specifies the operational components that instantiate the pre-registered protocol
$(\pi,\mathcal{C},\mathcal{A})$ (Section~\ref{sec:problem-formulation}). The emphasis is
\emph{controllability and auditability}: all decision rules, thresholds, denominators, and edge-case
behaviors are explicitly specified and frozen by a staged protocol, with Pilot locking the skeleton,
PreScreen locking admission-related values, and Calibration locking scene-routing values that require
coverage evidence.

\subsection{Meta-label collection and auditable missingness}

For each assignment $(w,t)$, the annotation UI requires a complete meta-label record $\mathbf{m}_{w,t}$
(cf.~\eqref{eq:meta}). In the actual Label Studio implementation, \texttt{difficulty} and \texttt{model\_issue}
are collected as required multi-select fields with UI-level mutual-exclusion rules (rather than separate
yes/no widgets for each tag). Any system-side missing field is recorded as NA and
reported as a process/format disclosure (not as a performance improvement).

To improve \emph{controllability} of Type~4 failures, we treat the meta-label requirements as part of the
collection protocol (not post-hoc cleaning freedom) and implement a two-layer ``hard validation'' defense.
First, the UI enforces ``at least one'' selection via required fields (e.g., \texttt{difficulty} is required;
\texttt{model\_issue} is required when collected). Second, we enforce mutual exclusion at submission time
to prevent ambiguous co-selection between a ``none/OK'' tag and any specific issue tag (\texttt{trivial} vs.
other difficulty factors; \texttt{acceptable} vs. other model failure modes): violating submissions are
blocked and recorded as auditable rejection events.

Because NA can still occur due to system/export issues (e.g., a field not transmitted, corrupted records, or
version mismatch), we do \emph{not} claim Type~4 is strictly zero. Instead, we disclose (i) NA rates with
explicit denominators and example tasks, and (ii) rejection counts and reasons (e.g., empty vs.
mutual-exclusion conflicts) as process evidence.

\subsection{PreScreen misleading initializations: source hierarchy and non-overlap}

\texttt{PreScreen\_manual} and \texttt{PreScreen\_semi} use disjoint image pools. The role of
\texttt{PreScreen\_manual} is to estimate baseline annotation ability without model initialization; hence its
expert-anchor subset should cover the major difficulty strata, but it does not need to mirror the failure-mode
composition of semi-automatic traps.

By contrast, \texttt{PreScreen\_semi} is designed to probe correction ability and blind trust under misleading
initialization. Its misleading initializations follow an asymmetric dual-source design: the primary source is a
pre-registered library of synthetic perturbation operators, while a small, fixed supplement may come from a
natural-failure trap bank curated from frozen model outputs on an independent candidate pool. The same
\texttt{base\_task\_id} must not be reused across manual and semi pools. If natural failures are insufficient for
some failure family, the realized quota for that family is reduced and the remaining semi-trap budget is filled
by the corresponding synthetic operator family; planned and realized source quotas are disclosed.

Clearly judgeable OOS or insufficient-evidence tasks that do not admit a stable geometric reference layout are
not used as main geometric traps. In \texttt{PreScreen\_manual}, such tasks may appear only as a fixed small
scope-gate subset scored on scope decisions rather than geometric IoU. In \texttt{PreScreen\_semi}, only an even
smaller OOS stress subset is allowed for blind-trust audit, and it is reported separately from the main geometric
trap pool. OOS cases whose scope boundary itself remains highly ambiguous are kept only in an ambiguity audit bank
rather than promoted to hard traps.

\subsection{Multi-label tag consensus (basic and weighted)}
\label{sec:tag-consensus}

We compute consensus separately for each binary tag $k$. Let $z_{t,w,k}\in\{0,1,\text{NA}\}$ denote worker
$w$'s response on task $t$. The operator outputs per-tag consensus and diagnostic uncertainty. The
consensus threshold $\tau_c$ and all edge-case rules are pre-registered.

\begin{algorithm}[t]
\caption{Per-tag (unweighted) consensus for multi-select meta-labels}
\label{alg:tag-consensus}
\KwRequire{Responses \ensuremath{\{z_{t,w,k}\}}, where \ensuremath{z_{t,w,k}\in\{0,1,\mathrm{NA}\}}; pre-registered threshold \ensuremath{\tau_c\in(0,1)}}
\KwEnsure{\ensuremath{c_{t,k}}, \ensuremath{\mathrm{conf}_{t,k}}, \ensuremath{\mathrm{disagree}_{t,k}} for each tag \ensuremath{k}; task-level set \ensuremath{C_t^{\mathrm{tag}}}; auditable NA rates}

\ForEach{$k\in\mathcal{K}$}{
	$n_{\mathrm{sel}}\leftarrow\sum_{w\in\mathcal{W}_t}\mathbb{I}[z_{t,w,k}=1]$\tcp{selected count}
	$n_{\mathrm{tot}}\leftarrow\sum_{w\in\mathcal{W}_t}\mathbb{I}[z_{t,w,k}\in\{0,1\}]$\tcp{valid count; NA excluded}
	\eIf{$n_{\mathrm{tot}}=0$}{
		$c_{t,k}\leftarrow\text{NA}$; $\mathrm{conf}_{t,k}\leftarrow\text{NA}$; $\mathrm{disagree}_{t,k}\leftarrow\text{NA}$\tcp{undefined; disclose}
	}{
		$p_{t,k}\leftarrow n_{\mathrm{sel}}/n_{\mathrm{tot}}$\tcp{vote rate}
		$c_{t,k}\leftarrow \mathbb{I}[p_{t,k}\ge\tau_c]$\tcp{tie handled by $\ge$}
		$\mathrm{conf}_{t,k}\leftarrow \max(p_{t,k},1-p_{t,k})$\tcp{margin confidence}
		$\mathrm{disagree}_{t,k}\leftarrow 1-\mathrm{conf}_{t,k}$\;
	}
}
$C_t^{\mathrm{tag}}\leftarrow\{k: c_{t,k}=1\}$\tcp{task-level tag set}
\Return all outputs and NA counts for audit reporting\;
\end{algorithm}

\paragraph{Weighted variant (controlled intervention).}
We additionally support a reliability-weighted consensus. To avoid circularity, weights must be computed
from PreScreen/Calibration only and must exclude the evaluated task (leave-one-task-out).
The main implementation records the event $\sum\omega_w=0$ as NA and reports its frequency; it does not
silently fall back to an unweighted vote in the main analysis.

The weighting cap is selected on the manual-anchor subset of \texttt{PreScreen\_manual}, which contains
20--22 expert-referenced anchors inside a 30-image pool. The main lock-in scheme is a repeated balanced
two-fold cross-fitting protocol; a three-fold stratified split is reported only as a sensitivity analysis.
Given the current anchor budget, five-fold cross-validation is not used in the main analysis because each
validation fold would be too thin to support stable family coverage.

\begin{algorithm}[t]
\caption{Per-tag (weighted) consensus (pre-registered edge-case behavior)}
\label{alg:weighted-tag-consensus}
\KwRequire{Responses \ensuremath{\{z_{t,w,k}\}}; weights \ensuremath{\{\omega_w\}} with \ensuremath{\omega_w\in[0,1]}; threshold \ensuremath{\tau_c}}
\KwEnsure{\ensuremath{c^{(\omega)}_{t,k}}, \ensuremath{\mathrm{conf}^{(\omega)}_{t,k}}, \ensuremath{\mathrm{disagree}^{(\omega)}_{t,k}}; auditable \ensuremath{W_{\mathrm{tot}}=0} rate}

\ForEach{$k\in\mathcal{K}$}{
	$W_{\mathrm{sel}}\leftarrow\sum_{w\in\mathcal{W}_t}\omega_w\,\mathbb{I}[z_{t,w,k}=1]$\tcp{weighted selected}
	$W_{\mathrm{tot}}\leftarrow\sum_{w\in\mathcal{W}_t}\omega_w\,\mathbb{I}[z_{t,w,k}\in\{0,1\}]$\tcp{weighted valid; NA excluded}
	\eIf{$W_{\mathrm{tot}}=0$}{
		$c^{(\omega)}_{t,k}\leftarrow\text{NA}$; $\mathrm{conf}^{(\omega)}_{t,k}\leftarrow\text{NA}$; $\mathrm{disagree}^{(\omega)}_{t,k}\leftarrow\text{NA}$\tcp{disclose; no silent fallback}
	}{
		$p^{(\omega)}_{t,k}\leftarrow W_{\mathrm{sel}}/W_{\mathrm{tot}}$\tcp{weighted vote rate}
		$c^{(\omega)}_{t,k}\leftarrow \mathbb{I}[p^{(\omega)}_{t,k}\ge\tau_c]$\;
		$\mathrm{conf}^{(\omega)}_{t,k}\leftarrow \max(p^{(\omega)}_{t,k},1-p^{(\omega)}_{t,k})$\;
		$\mathrm{disagree}^{(\omega)}_{t,k}\leftarrow 1-\mathrm{conf}^{(\omega)}_{t,k}$\;
	}
}
\Return all outputs and the frequency of $W_{\mathrm{tot}}=0$\;
\end{algorithm}

\subsection{Quality and reliability estimation (leave-one-out protocol)}
\label{sec:loo-reliability}

We separate \emph{quality} (task-level agreement) from \emph{reliability} (worker-level robustness).
To avoid self-inflation, worker evaluation follows a leave-one-out (LOO) protocol.

\begin{algorithm}[t]
\caption{LOO reliability estimation on multi-annotated tasks}
\label{alg:loo-reliability}
\KwRequire{A set of tasks \ensuremath{\mathcal{T}_{\mathrm{cal}}} with redundancy \ensuremath{k\ge k_{\min}}; submissions \ensuremath{\{A_{t,w}\}}; IoU metric; pre-registered \ensuremath{k_{\min}} and reporting policy for \ensuremath{k=3} tasks}
\KwEnsure{Task agreement proxy \ensuremath{\mathrm{IAA}_t}; LOO proxy \ensuremath{\mathrm{IoU}_{\mathrm{LOO}}(t,w)}; worker reliability \ensuremath{r_w}; disclosed rates of degeneracy/missingness}

\ForEach{$t\in\mathcal{T}_{\mathrm{cal}}$}{
	\If{$k(t)<k_{\min}$}{
		Mark task $t$ as \texttt{degenerate}; disclose its frequency\tcp{no silent pruning}
		\textbf{continue}\;
	}
	$\mathrm{IAA}_t\leftarrow \mathrm{median}_{w\ne w'}\;\mathrm{IoU}(A_{t,w},A_{t,w'})$\tcp{task agreement}
	$C_t\leftarrow \arg\max_{w}\ \mathrm{median}_{w'\ne w}\ \mathrm{IoU}(A_{t,w},A_{t,w'})$\tcp{consensus representative}
	\ForEach{$w\in\mathcal{W}_t$}{
		$C_t^{(-w)}\leftarrow \arg\max_{u\in\mathcal{W}_t\setminus\{w\}}\ \mathrm{median}_{v\in\mathcal{W}_t\setminus\{w,u\}}\ \mathrm{IoU}(A_{t,u},A_{t,v})$\tcp{LOO representative}
		$\mathrm{IoU}_{\mathrm{LOO}}(t,w)\leftarrow \mathrm{IoU}(A_{t,w},C_t^{(-w)})$\tcp{worker-on-task proxy}
	}
}
\ForEach{$w\in\mathcal{W}$}{
	$r_w\leftarrow \mathrm{median}_{t\in\mathcal{T}_w}\ \mathrm{IoU}_{\mathrm{LOO}}(t,w)$\tcp{global reliability}
	Compute a pre-registered bootstrap CI and record its half-width\tcp{stopping criterion}
}
\Return all estimates and disclosure rates\;
\end{algorithm}

\subsection{Worker profiling: reliability axis and discrete risk tiers}

The main analysis does not use a single continuous ``spammer score.'' Instead, worker profiling follows a
\emph{reliability axis + discrete risk tier} design that is easier to audit under limited sample size.
The vertical axis is $\mathrm{LCB}(r_w)$, while the horizontal axis is a discrete risk tier
$R_w^{\mathrm{tier}}\in\{R0,R1,R2,R3\}$.

The tier is determined by auditable process evidence rather than a weighted sum of heterogeneous proxies.
We use three core risk signatures: blind-trust risk $T_w$ from misleading-initialization semi tasks,
condition fragility $C_w$ from high-risk bucket degradation, and gate-failure tendency $G_w$ from
computable engineering failures. Low-effort proxies and NA/missingness are retained only as auxiliary audit
fields rather than primary profile axes.

The tiers are assigned in a fixed order: $R3$ (unusable/noisy) if reliability is below the conservative
floor or gate-failure tendency is too high; $R2$ (condition-fragile) if high-risk buckets fail despite an
acceptable global level; $R1$ (trust-vulnerable) if blind trust remains high without bucket failure; and
$R0$ (stable) otherwise. This priority rule is exported together with a machine-readable reason code so that
the profile remains reproducible.

\subsection{Risk-aware task distribution (stress mode and sequential redundancy)}
\label{sec:routing-policy}

Routing is treated as a budgeted quality-control policy: higher-risk tasks are prioritized to more stable
workers, and redundancy is spent where it reduces disagreement. We distinguish (i) pre-annotation risk
proxies $(d_t,g_t)$ available before labels, and (ii) post-submission diagnostics such as tag disagreement.

In operational terms, high-risk tasks reserve the primary pool for $R0$ workers whose
$\mathrm{LCB}(r_w)$ exceeds the conservative threshold. $R2$ workers are excluded from the strata where they
have already failed, $R1$ workers are not used as the primary pool for misleading-initialization-heavy semi
tasks, and $R3$ workers are excluded from the main routing pool and from weighted consensus.

\begin{algorithm}[t]
\caption{Risk-aware routing with stress mode and stopping rule}
\label{alg:routing}
\KwRequire{Tasks \ensuremath{\mathcal{T}} with risk proxies \ensuremath{d_t,g_t}; worker pool with reliability lower bounds \ensuremath{\mathrm{LCB}(r_w)} (and optional strata LCBs); budget \ensuremath{B}; pre-registered thresholds; max redundancy \ensuremath{k_{\max}}}
\KwEnsure{Assignments \ensuremath{\pi(t)} and auditable routing decisions under a fixed budget}

\ForEach{$t\in\mathcal{T}$}{
	Determine risk level using $(d_t,g_t)$ and pre-registered triggers\tcp{cold-start proxies}
	\eIf{stress mode is triggered (e.g., $d_t\ge\tau_{\mathrm{stress}}$ or $g_t$ flagged)}{
		Candidate workers $\mathcal{W}_t\leftarrow\{w: \mathrm{LCB}(r_w)\ge\tau_{\mathrm{high}}\}$\tcp{conservative pool}
		Initialize redundancy $k\leftarrow 3$\tcp{pre-registered}
	}{
		Candidate workers $\mathcal{W}_t\leftarrow\{w: \mathrm{LCB}(r_w)\ge\tau_{\mathrm{low}}\}$\tcp{default pool}
		Initialize redundancy $k\leftarrow k_0(t)$\tcp{risk-stratified start}
	}

	Assign $k$ workers from $\mathcal{W}_t$ under a load constraint and remaining budget\tcp{no domination}
	Collect submissions and meta-labels; compute tag disagreement via Algorithm~\ref{alg:tag-consensus}\;

	\While{disagreement is above threshold \textbf{and} $k<k_{\max}$ \textbf{and} budget remains}{
		Assign one additional worker from $\mathcal{W}_t$\tcp{sequential add-1}
		Update disagreement and audit records\;
		$k\leftarrow k+1$\;
	}

	\If{$k=k_{\max}$ and disagreement remains high}{
		Mark task $t$ for expert review and disclose its rate\tcp{auditable fallback}
	}
}
\Return $\pi(\cdot)$, redundancy used per task, and all disclosure rates\;
\end{algorithm}

Sensitivity checks (e.g., varying $\tau_c$, $k_{\max}$, or the weight mapping cap) are reported as
pre-registered grids and are not used for post-hoc selection.
